\chapter{Hamiltonian jets and quantum currents}
\label{chap:hamiltonian-jets}

A finite Taylor expansion can carry a nontrivial Hamiltonian symmetry
when the symmetry fixes the expansion point. Translations violate that
condition: differentiation can return a discarded coefficient to the
retained range. Removing translations therefore changes the representation
problem, and it permits finite matter modules with an exact quantum action.

The first module is spanned by $x,y$. Quadratic Hamiltonians act by
$\mathfrak{sl}_2$, yet their symmetric cubic trace is zero.
The same cancellation persists at every Taylor order. Higher Hamiltonians
only raise matter degree, so a closed matrix word sees only the quadratic
blocks. A bilinear form on each block reverses that word with a sign.
Together with the odd propagator and background signs, this identity
cancels every complete scalar wheel in the free theory on $\mathbb C^2$.

We construct the resulting current maps and the unital BV algebras they
generate inside the actual distribution observables. Their inverse limit
retains multiplication, the differential, the BV operator, and every higher
current component. Although individual currents have zero BV Laplacian,
products can have a nonzero contraction. A removed pair of matter modes
shows why the full ambient BV algebras do not admit the same projections.
Finally, Taylor expansion of a spatial Hamiltonian supplies the transport
term needed to compare internal symmetry with motion of insertion points.

\section{A finite Taylor module that retains the action}

Work over \(\mathbb C\). The variables \(x,y\) in this section are formal internal variables. Put
\[
R=\mathbb C[[x,y]],\qquad \mathfrak m=(x,y),\qquad
H_+=\mathfrak m^2,
\]
with
\[
\{f,g\}=f_xg_y-f_yg_x.
\]
We identify \(\mathfrak m^2\) with its image in \(R/\mathbb C\).
Give the coefficient field the discrete topology for this formal
construction. The subspaces $\mathfrak m^r$ define its neighborhoods
of zero. All Lie algebras and matter modules in this section are
concentrated in cohomological degree zero.

For \(N\geq1\), define
\begin{equation}
V_N=\mathfrak m/\mathfrak m^{N+1},
\qquad
\mathfrak g_N=\mathfrak m^2/\mathfrak m^{N+2}.
\label{eq:hj-1}
\end{equation}
The action is
\begin{equation}
\rho_N(\overline f)(\overline g)
   =\overline{\{f,g\}}\in V_N.
\label{eq:hj-2}
\end{equation}

\begin{proposition}\label{prop:hj-faithful}
For every $N\geq1$, $\mathfrak g_N$ is a Lie algebra and
\eqref{eq:hj-2} is a faithful representation on $V_N$.
\end{proposition}
\begin{proof}
For \(a,b\geq1\),
\begin{equation}
\{\mathfrak m^a,\mathfrak m^b\}
   \subset\mathfrak m^{a+b-2}.
\label{eq:hj-3}
\end{equation}
Consequently:

\begin{itemize}
\item \(\mathfrak m^{N+2}\) is a Lie ideal of \(\mathfrak m^2\).
\item \(\mathfrak m^{N+1}\) is preserved by \(\mathfrak m^2\).
\item Changing \(f\) by \(\mathfrak m^{N+2}\) changes \(\{f,g\}\), for \(g\in\mathfrak m\), by \(\mathfrak m^{N+1}\).
\end{itemize}

Thus both quotients and the action are well-defined. Jacobi gives
\begin{equation}
[\rho_N(f),\rho_N(g)]=\rho_N(\{f,g\}).
\label{eq:hj-4}
\end{equation}

The representation is faithful. If \(\rho_N(f)=0\), evaluating on \(x,y\) gives
\[
-f_y,\ f_x\in\mathfrak m^{N+1}.
\]
All homogeneous components of \(f\) of degrees \(2,\ldots,N+1\) therefore vanish. Characteristic zero is used here. Hence \(f=0\) in \(\mathfrak g_N\).
\end{proof}

Writing \(H_d=\mathbb C[x,y]_d\), we have
\[
V_N=\bigoplus_{d=1}^N H_d,
\qquad
\mathfrak g_N=\bigoplus_{r=2}^{N+1}H_r,
\]
and
\begin{equation}
\dim V_N=\frac{N(N+3)}2,\qquad
\dim\mathfrak g_N=\frac{N(N+5)}2.
\label{eq:hj-5}
\end{equation}

The projection \(H_+\to\mathfrak g_N\) is continuous for the formal topology because its kernel is open. Its composite with \(\rho_N\) is a continuous finite-dimensional representation. This remains true when the finite-dimensional target has its usual complex topology.
Write $\pi_N:H_+\to\mathfrak g_N$ for this quotient map.

If \(f_r\in H_r\), then
\begin{equation}
\rho_N(f_r):H_d\longrightarrow H_{d+r-2},
\label{eq:hj-6}
\end{equation}
with the target interpreted as zero when \(d+r-2>N\). Thus every action matrix is nondecreasing in internal degree.

\section{The first two jets}

The first quotient contains the full quadratic action. The next one adds a degree-raising block, which reveals why higher Taylor terms cannot close an internal loop.

Put
\[
e=x^2/2,\qquad f=-y^2/2,\qquad h=-xy.
\]
Direct Poisson brackets give
\begin{equation}
[h,e]=2e,\qquad[h,f]=-2f,\qquad[e,f]=h.
\label{eq:hj-8}
\end{equation}
These are the defining brackets of $\mathfrak{sl}_2$.
In matrix expressions, a Hamiltonian denotes its action on the
specified module. Columns record the images of basis vectors.

For \(N=1\), the ordered basis \(x,y\) gives
\[
h=\begin{pmatrix}1&0\\0&-1\end{pmatrix},\qquad
e=\begin{pmatrix}0&1\\0&0\end{pmatrix},\qquad
f=\begin{pmatrix}0&0\\1&0\end{pmatrix}.
\]
Thus
\begin{equation}
\operatorname{Tr}(h^2)=2,\quad
\operatorname{Tr}(h^3)=0,\quad
\operatorname{Tr}(h^4)=2,\quad
\operatorname{Tr}(hef)=1.
\label{eq:hj-24}
\end{equation}

For \(N=2\), use \(x,y,x^2,xy,y^2\). Then
\[
h=\operatorname{diag}(1,-1,2,0,-2),
\]
and the other two action matrices are
\[
e=\begin{pmatrix}
0&1&0&0&0\\
0&0&0&0&0\\
0&0&0&1&0\\
0&0&0&0&2\\
0&0&0&0&0
\end{pmatrix},\qquad
f=\begin{pmatrix}
0&0&0&0&0\\
1&0&0&0&0\\
0&0&0&0&0\\
0&0&2&0&0\\
0&0&0&1&0
\end{pmatrix}.
\]
Their traces give
\begin{equation}
\operatorname{Tr}(h^2)=10,\quad
\operatorname{Tr}(h^3)=0,\quad
\operatorname{Tr}(h^4)=34,\quad
\operatorname{Tr}(hef)=5.
\label{eq:hj-25}
\end{equation}
The four cubic Hamiltonians act by
\[
\begin{array}{c|cc}
 &x&y\\ \hline
x^3&0&3x^2\\
x^2y&-x^2&2xy\\
xy^2&-2xy&y^2\\
y^3&-3y^2&0
\end{array}
\]
and annihilate \(H_2\). Their matrices are strictly degree-raising, so none contributes to a trace word.
For example, with the decomposition $V_2=H_1\oplus H_2$,
\[
\rho_2(x^3)=
\begin{pmatrix}0_{2\times2}&0_{2\times3}\\
C_{x^3}&0_{3\times3}\end{pmatrix},
\qquad
C_{x^3}=\begin{pmatrix}0&3\\0&0\\0&0\end{pmatrix}.
\]
The remaining lower-left blocks, in the same bases, are
\[
C_{x^2y}=\begin{pmatrix}-1&0\\0&2\\0&0\end{pmatrix},\quad
C_{xy^2}=\begin{pmatrix}0&0\\-2&0\\0&1\end{pmatrix},\quad
C_{y^3}=\begin{pmatrix}0&0\\0&0\\-3&0\end{pmatrix}.
\]
The new matter degrees affect traces of quadratic Hamiltonians.
The new Hamiltonian degrees cannot return an input to its original
degree. This distinction persists at every $N$.

\section{Why only quadratic Hamiltonians enter a trace}

A closed internal contraction is a trace of an action word.
Its diagonal coefficients must return to the starting matter degree.
Let \(f_{j,2}\) denote the quadratic component of \(f_j\in H_+\).
In a trace over $H_d$, $\rho_d(q)$ means its restriction to the
invariant summand $H_d\subset V_d$, for a quadratic Hamiltonian $q$.
Equation~\eqref{eq:hj-6} gives, for every word length \(r\geq1\),
\begin{equation}
\boxed{
\begin{aligned}
&\operatorname{Tr}_{V_N}
 \bigl(\rho_N(f_1)\cdots\rho_N(f_r)\bigr)
 \\ &\qquad={}
 \sum_{d=1}^N
 \operatorname{Tr}_{H_d}
 \bigl(\rho_d(f_{1,2})\cdots\rho_d(f_{r,2})\bigr).
\end{aligned}
}
\label{eq:hj-7}
\end{equation}
Indeed, each factor raises internal degree by a nonnegative amount. A diagonal block of a product can therefore use only the degree-preserving part of every factor. These are precisely the quadratic Hamiltonian parts.

In particular, any trace word containing one Hamiltonian with zero quadratic component vanishes. This includes all higher Taylor components, without any genericity assumption.

On
\[
v_a=x^ay^{d-a},\qquad 0\leq a\leq d,
\]
their actions are
\begin{equation}
hv_a=(2a-d)v_a,\qquad
ev_a=(d-a)v_{a+1},\qquad
fv_a=av_{a-1}.
\label{eq:hj-9}
\end{equation}
The product rule identifies this action with the $d$th symmetric
power of the action on $\mathbb Cx\oplus\mathbb Cy$.

Define a bilinear form on \(H_d\) by
\begin{equation}
B_d(v_a,v_b)=0\quad(a+b\ne d),
\qquad
B_d(v_a,v_{d-a})=\frac{(-1)^a}{\binom da}.
\label{eq:hj-10}
\end{equation}
It is nondegenerate. Formula \eqref{eq:hj-9} directly verifies
\begin{equation}
B_d(Xv,w)+B_d(v,Xw)=0
\qquad(X=e,f,h).
\label{eq:hj-11}
\end{equation}
For example, invariance under \(e\) reduces to
\[
(d-a)B_d(v_{a+1},v_{d-a-1})
 +(a+1)B_d(v_a,v_{d-a})=0.
\]
The displayed coefficients in \eqref{eq:hj-10} satisfy this identity. The checks for \(f,h\) are immediate from \eqref{eq:hj-9}.

If \(J_d\) is the matrix of \(B_d\), every quadratic action matrix satisfies
\[
X^{\mathsf T}J_d+J_dX=0.
\]
Transposing a product and conjugating by \(J_d\) gives
\[
\operatorname{Tr}_{H_d}(X_1\cdots X_r)
 =(-1)^r\operatorname{Tr}_{H_d}(X_r\cdots X_1).
\]
Combining this with \eqref{eq:hj-7} proves the full jet identity
\begin{equation}
\boxed{
\begin{aligned}
&\operatorname{Tr}_{V_N}
 \bigl(\rho_N(f_1)\cdots\rho_N(f_r)\bigr)
 \\ &\qquad=(-1)^r
\operatorname{Tr}_{V_N}
 \bigl(\rho_N(f_r)\cdots\rho_N(f_1)\bigr).
\end{aligned}
}
\label{eq:hj-12}
\end{equation}
The full jet representation need not preserve the forms \eqref{eq:hj-10}. The trace identity follows from its degree-preserving blocks.

\section{The cubic invariant and ordered color factors}

The local scalar curvature in complex dimension two uses a symmetric
cubic trace. An individual orientation of a loop instead has an ordered
matrix word. To distinguish them, define
\begin{equation}
T_N(f_1,f_2,f_3)
 =\frac16\sum_{\sigma\in S_3}
 \operatorname{Tr}_{V_N}
 \bigl(\rho_N(f_{\sigma(1)})
       \rho_N(f_{\sigma(2)})
       \rho_N(f_{\sigma(3)})\bigr).
\label{eq:hj-13}
\end{equation}
\begin{proposition}\label{prop:hj-cubic}
The complete symmetric cubic tensor of $\rho_N$ vanishes:
\begin{equation}
\boxed{T_N=0\qquad\text{for every }N\geq1.}
\label{eq:hj-14}
\end{equation}
\end{proposition}
\begin{proof}
Reversal pairs each summand of \eqref{eq:hj-13} with its negative
by \eqref{eq:hj-12}. This applies to arbitrary Hamiltonian inputs,
including every higher Taylor component.
\end{proof}

Individual ordered cubic traces can be nonzero. Cyclicity and \eqref{eq:hj-12} give
\begin{equation}
\operatorname{Tr}(XYZ)
 =-\operatorname{Tr}(XZY)
 =\frac12\operatorname{Tr}\bigl(X[Y,Z]\bigr).
\label{eq:hj-15}
\end{equation}
This also computes the entire ordered cubic tensor in terms of the quadratic trace pairing.

The vanishing in \eqref{eq:hj-14} is thus stronger than vanishing on a selected Cartan element, but weaker than vanishing of every ordered cubic word.

\section{Cartan moments as the jet order increases}

The cubic invariant is already zero, but a limit of ordinary traces
would also have to accommodate the even moments. They provide a separate
test of what passage to increasing Taylor order can preserve.
The eigenvalues of \(\rho_N(h)\) are
\[
2a-d,\qquad 1\leq d\leq N,\quad 0\leq a\leq d.
\]
Hence
\begin{equation}
\operatorname{Tr}_{V_N}(h^r)
 =\sum_{d=1}^N\sum_{a=0}^d(2a-d)^r.
\label{eq:hj-16}
\end{equation}
All odd moments vanish by pairing \(a\) with \(d-a\).

For a single degree \(d\),
\begin{equation}
\sum_{a=0}^d(2a-d)^2
 =\frac{d(d+1)(d+2)}3,
\label{eq:hj-17}
\end{equation}
and
\begin{equation}
\sum_{a=0}^d(2a-d)^4
 =\frac{d(d+1)(d+2)(3d^2+6d-4)}{15}.
\label{eq:hj-18}
\end{equation}
One direct derivation expands
\[
\sum_{a=0}^d e^{(2a-d)t}
 =\frac{\sinh((d+1)t)}{\sinh t}
\]
as a formal power series through \(t^4\).
More explicitly, this quotient is
\[
(d+1)\left(
1+\frac{(d+1)^2-1}{6}t^2
 +\frac{3(d+1)^4-10(d+1)^2+7}{360}t^4
 +O(t^6)\right).
\]
Multiplication of the two coefficients by $2!$ and $4!$ gives
\eqref{eq:hj-17} and \eqref{eq:hj-18}.

Summing \eqref{eq:hj-17}--\eqref{eq:hj-18} gives
\begin{equation}
\boxed{
\operatorname{Tr}_{V_N}(h^2)
 =\frac{N(N+1)(N+2)(N+3)}{12},
}
\label{eq:hj-19}
\end{equation}
\begin{equation}
\boxed{\operatorname{Tr}_{V_N}(h^3)=0,}
\label{eq:hj-20}
\end{equation}
\begin{equation}
\boxed{
\operatorname{Tr}_{V_N}(h^4)
 =
\frac{N(N+1)(N+2)(N+3)(2N^2+6N-3)}{60}.
}
\label{eq:hj-21}
\end{equation}
The sums can also be verified without any summation formula: the difference between the right sides at \(N\) and \(N-1\) is respectively \eqref{eq:hj-17} and \eqref{eq:hj-18}, and both vanish at \(N=0\).

For an arbitrary quadratic
\[
q=ae+bh+cf,
\]
its matrix on \(H_1\) is
\[
\begin{pmatrix}b&a\\ c&-b\end{pmatrix}.
\]
Consequently,
\begin{equation}
\operatorname{Tr}_{V_N}\rho_N(q)^2
 =\operatorname{Tr}_{V_N}(h^2)(b^2+ac),
\label{eq:hj-22}
\end{equation}
\begin{equation}
\operatorname{Tr}_{V_N}\rho_N(q)^4
 =\operatorname{Tr}_{V_N}(h^4)(b^2+ac)^2.
\label{eq:hj-23}
\end{equation}
To prove these formulas, first diagonalize the displayed matrix when \(b^2+ac\ne0\). Equation \eqref{eq:hj-9} then determines its action on every symmetric power. Both sides are polynomial in \(a,b,c\), so the identities extend to \(b^2+ac=0\).

\section{Internal Hamiltonians and compact spatial sources}

To form currents, place these finite internal modules over a spatial
manifold. A matrix acts on the matter coefficient at an insertion point;
the propagator connects different insertion points. The spatial manifold is
\[
X=\mathbb C^2_z.
\]
Its coordinates \(z_1,z_2\) are distinct from the internal formal variables \(x,y\).

For each \(N\), define the differential graded Lie algebra
\[
\mathfrak a_N=\Omega_c^{0,\bullet}(X,\mathfrak g_N),
\qquad
d=\bar\partial,
\]
\begin{equation}
[\alpha\otimes f,\eta\otimes g]
 =(\alpha\wedge\eta)\otimes\{f,g\}.
\label{eq:hj-26}
\end{equation}
These smooth coefficients use the usual complex topology. On each fixed compact support, use the ordinary smooth seminorms.
Thus, for a compact $K\subset X$, a seminorm bounds the coefficients
and their real spatial derivatives through an assigned order on $K$.
The bracket is well-defined because $\pi_N$ is a Lie quotient, and
$\bar\partial$ is a derivation of the wedge product.
The quotient maps $\mathfrak g_{N+1}\to\mathfrak g_N$ induce
surjective differential graded Lie maps on these sources.

A minimal formal Hamiltonian current source is the algebraic differential graded Lie algebra
\begin{equation}
\Omega_c^{0,\bullet}(X)\otimes_{\mathbb C}H_+.
\label{eq:hj-27}
\end{equation}
Its maps
\begin{equation}
q_N=\operatorname{id}\otimes\pi_N:
\Omega_c^{0,\bullet}(X)\otimes H_+
 \longrightarrow\mathfrak a_N
\label{eq:hj-28}
\end{equation}
commute with the differential and bracket. Every construction below can be pulled back along \eqref{eq:hj-28}. Its continuity statements hold on the specified finite families and fixed compact supports.

For the inverse construction, an explicit larger source is
\begin{equation}
\mathfrak a_+=\varprojlim_N\mathfrak a_N.
\label{eq:hj-29}
\end{equation}
An element has smooth formal Taylor coefficients such that each finite internal jet has compact spatial support. A common compact support for all coefficients is allowed but is not imposed by \eqref{eq:hj-29}. Formula \eqref{eq:hj-27} embeds in \eqref{eq:hj-29}.
The inverse limit is degreewise. Its coordinate projections are the
maps to $\mathfrak a_N$; no interchange of a tensor product with this
inverse limit is assumed. Fixed-support continuity below always refers
to the ordinary seminorms at each finite $N$.

This definition avoids identifying smooth formal coefficients with smooth maps into a vector space whose coefficient field was given the discrete topology.

\section{The finite matter theory and its path kernels}

The matrix action gives the local coupling $\int\beta\rho_N(A)\gamma\Omega$.
Its Gaussian contractions produce paths with two external matter legs,
and cycles with none. Fix \(N\) and a positive heat scale \(L\).
Take matter fields
\[
\gamma\in\Omega^{0,\bullet}(X,V_N),
\qquad
\beta\in\Omega^{0,\bullet}(X)\otimes V_N^\vee[1].
\]
Use
\[
\Omega=dz_1\wedge dz_2,\qquad
d\bar z_1d\bar z_2\Omega=4\,dV_4.
\]
The total degrees of \(\beta,A,\gamma\) are \(1,1,0\). Ghost coefficients precede forms:
\begin{equation}
(u\alpha)(v\eta)
 =(-1)^{|\alpha|\operatorname{gh}(v)}
   uv\,\alpha\wedge\eta.
\label{eq:hj-30}
\end{equation}
The shift convention is $(E[k])^j=E^{j+k}$.
A form-degree-$q$ component of $\gamma,\beta,A$ has ghost number
$-q,1-q,1-q$, respectively. Here $\beta$ is the coefficient field:
its full field is $\beta\wedge\Omega$.
If $z_j=s_j+ir_j$, then
$dV_4=ds_1\wedge dr_1\wedge ds_2\wedge dr_2$.

Evaluate the universal background on any finite family
\[
A=\sum_i t_i\alpha_i,\qquad
\alpha_i\in\mathfrak a_N^{p_i},\qquad |t_i|=1-p_i.
\]
The parameters form a free graded commutative algebra.
Odd parameters square to zero. Identities mean identities of every
ordered parameter coefficient. One may adjoin the finitely many
differentials and brackets occurring in that coefficient.
No finite-dimensional background subspace is assumed closed under them.
The scalar $\hbar$ is a formal parameter of degree zero.

Write
\[
M_N(A)=\rho_N(A).
\]
For \(0<\epsilon<L\), use the fixed propagator
\begin{equation}
P_{\epsilon,L}(u)
 =\int_\epsilon^L
   \frac{e^{-|u|^2/(4t)}}{64\pi^2t^3}
   \bigl(\bar u_1d\bar u_2-\bar u_2d\bar u_1\bigr)\,dt.
\label{eq:hj-31}
\end{equation}
All forms use full differences. In particular,
\begin{equation}
P_{\epsilon,L}(-u)=P_{\epsilon,L}(u).
\label{eq:hj-32}
\end{equation}
For $u=z-w$, this means $d\bar u_j=d\bar z_j-d\bar w_j$.
The corresponding heat form is
\[
H_t(u)=\frac{e^{-|u|^2/(4t)}}{64\pi^2t^2}
              d\bar u_1\wedge d\bar u_2,
\qquad dP_{\epsilon,L}=H_\epsilon-H_L.
\]
The internal pairing is evaluation $V_N^\vee\otimes V_N\to\mathbb C$.

Let $E_N$ be the graded finite-rank bundle of matter components.
The algebra of observables is
\[
\mathcal O_{\mathrm{dist},N}
=\bigoplus_{m\geq0}
\mathcal E'_c\bigl(X^m;(E_N^\vee)^{\boxtimes m}
                     \otimes\operatorname{Dens}\bigr)_{S_m}.
\]
Its elements are compact distribution kernels, modulo graded symmetric
permutations of the matter slots. The summand $m=0$ is $\mathbb C$.
A scalar kernel $T$ has a bound
\[
|T(\phi)|\leq C\max_{|\nu|\leq m}\sup_K|\partial^\nu\phi|
\]
for some compact $K$ and finite $m$.
Its strong topology is uniform convergence on families of smooth tests
bounded in every derivative on each compact set.
Products keep position variables separate by external tensor product.
The free differential $Q$ acts by transposing the matter differential
onto tests. A positive-scale contraction inserts its smooth kernel into
a test and pushes forward over the contracted slots.
These are the operations constructed in
Chapter~\ref{chap:quantum-current-anomaly}.

For the paired heat tensor $K_t$, use the signs
\[
\Delta_t=-\partial_{K_t},\qquad
\mathcal C_{\epsilon,L}=\partial_{P_{\epsilon,L}},\qquad
[Q,\mathcal C_{\epsilon,L}]=\Delta_\epsilon-\Delta_L.
\]
Odd coordinate derivatives are left derivatives. The propagator form
stands on the left of its contraction operator.
The operator $\Delta_t$ and its derived bracket have degree one.
They satisfy $\Delta_t^2=0$, $Q\Delta_t+\Delta_tQ=0$, and the
graded Jacobi and derivation identities proved there.
The singular $P_{0,L}$ is used only in the path kernels below,
not as an operator on arbitrary distributions.

For
\[
s_n=(-1)^{(n-1)(n-2)/2},
\]
the tree interaction is
\begin{equation}
\begin{aligned}
T_{n,N}^{\epsilon,L}(A)
={}&s_n\int_{X^n}P_{12}\cdots P_{n-1,n}\\
&\quad\cdot
 \beta(z_1)M_N(A(z_1))\cdots
 M_N(A(z_n))\gamma(z_n)\,
 \Omega_1\cdots\Omega_n.
\end{aligned}
\label{eq:hj-33}
\end{equation}
For \(n=1\), the propagator product is \(1\).

These kernels have limits \(T_{n,N}^L\). Indeed,
\[
\|P_{0,L}\|_{L^1}\leq C\sqrt L,\qquad
\|P_{0,L}-P_{\epsilon,L}\|_{L^1}
 \leq C\sqrt\epsilon.
\]
For $n\geq2$, the $n-1$ edge differences and one base point are
independent coordinates, with absolute real Jacobian one.
A telescoping product expansion therefore gives, on a fixed compact support,
\begin{equation}
\bigl|(T_{n,N}^{\epsilon,L}-T_{n,N}^L)(\Phi)\bigr|
 \leq C_{n,N,K}\sqrt\epsilon\,L^{(n-2)/2}
           \|\Phi\|_{C^0}.
\label{eq:hj-34}
\end{equation}
The resulting observable is quadratic, with a compact distribution kernel in its two matter slots.
In this estimate $\Phi$ includes the smooth background coefficients
and the smooth test in the two remaining matter slots.
On a bounded family of backgrounds with fixed support, its bound is
uniform. The independent edge variables cover partial collisions as well
as the full diagonal. For $n=1$, the kernel is independent of $\epsilon$.

The scalar wheel formula is
\begin{equation}
\begin{aligned}
W_{n,N}^{\epsilon,L}(A)
 ={}&\frac{s_n}{n}\int_{X^n} P_{12}\cdots P_{n1}\\
 &\quad\cdot
 \operatorname{Tr}_{V_N}
 \bigl(M_N(A(z_1))\cdots M_N(A(z_n))\bigr)
 \Omega_1\cdots\Omega_n.
\end{aligned}
\label{eq:hj-35}
\end{equation}
for \(n\geq3\), with \(W_{1,N}=W_{2,N}=0\).

The connected Gaussian expansion in Chapter~\ref{chap:all-arity-current-actions} gives the factor \(1/n\) and sign \(s_n\). A path uses all \(n!\) vertex orders, whereas a cycle identifies its \(n\) cyclic rotations. Specializing the representation does not change these counts.

\section{Reversal cancels the complete wheels}

An ordered trace can be nonzero, as $\operatorname{Tr}(hef)$ shows.
The quantum term sums both orientations with their form and ghost signs.
Those additional signs determine whether a scalar correction remains.
\begin{proposition}\label{prop:hj-wheel}
At every $0<\epsilon<L$, all complete scalar wheel polynomials for
$\rho_N$ vanish. Their limiting polynomials and the local cubic
curvature vanish as well.
\end{proposition}
\begin{proof}
Each \(A_i\) has total degree one. Expand each matrix-valued \(A_i\) into an odd scalar coefficient times an ordinary Hamiltonian matrix. Reversing the scalar coefficients contributes
\[
(-1)^{n(n-1)/2}.
\]
Reversing the ordinary matrix trace contributes \((-1)^n\) by \eqref{eq:hj-12}. Hence
\begin{equation}
\operatorname{Tr}(M_N(A_n)\cdots M_N(A_1))
 =
(-1)^{n(n+1)/2}
\operatorname{Tr}(M_N(A_1)\cdots M_N(A_n)).
\label{eq:hj-36}
\end{equation}

In \eqref{eq:hj-35}, reverse the vertex labels \(i\mapsto n+1-i\). The volume forms introduce no sign. Equation \eqref{eq:hj-32} fixes each reversed edge form. Reordering the first \(n-1\) odd edge forms back to their original order contributes
\[
(-1)^{(n-1)(n-2)/2}.
\]
Together with \eqref{eq:hj-36}, the total sign is
\[
(-1)^{n(n+1)/2+(n-1)(n-2)/2}
 =(-1)^{n^2-n+1}=-1.
\]
Thus
\[
W_{n,N}^{\epsilon,L}(A)=-W_{n,N}^{\epsilon,L}(A),
\]
and
\begin{equation}
\boxed{
W_{n,N}^{\epsilon,L}=W_{n,N}^{L}=0
\quad\text{for every }n.
}
\label{eq:hj-37}
\end{equation}

This cancellation holds for the complete wheel polynomial. It does not say that each ordered color trace or each labeled oriented graph is zero.

The cubic local curvature also vanishes directly. Its prescribed expression is
\begin{equation}
\Theta_N(A)
 =-\frac1{24\pi^2}\int_X
 \operatorname{Tr}_{V_N}
 \left[
 A(\partial_{z_1}A\,\partial_{z_2}A
    -\partial_{z_2}A\,\partial_{z_1}A)
 \right]\Omega.
\label{eq:hj-38}
\end{equation}
Here and below \(A\) in a matrix product means \(M_N(A)\).

All three matrix-valued factors in either cubic trace have total degree one. Formula \eqref{eq:hj-36} at \(n=3\), followed by graded cyclicity, identifies the two traces in \eqref{eq:hj-38}. Consequently,
\begin{equation}
\boxed{\Theta_N=0.}
\label{eq:hj-39}
\end{equation}
This agrees with the vanishing of the complete symmetric cubic tensor in \eqref{eq:hj-14}.
For $n=1$, the propagator has $P_{\epsilon,L}(0)=0$.
For $n=2$, its full-difference pullback gives the product
$P_{\epsilon,L}(u)^2=0$. These supply the two cases excluded
from the formula \eqref{eq:hj-35}.
\end{proof}

Define the formal interaction
\begin{equation}
I_{N,L}(A)=\sum_{n\geq1}T_{n,N}^L(A),
\label{eq:hj-40}
\end{equation}
completed in background weight. It is interpreted on finite parameter families
\[
A=\sum_i t_i\alpha_i,\qquad |t_i|=1-|\alpha_i|.
\]
At each background order, every expression is finite. No numerical convergence of \eqref{eq:hj-40} is asserted.

Theorem~\ref{thm:aac-fixed-scale}, with its algebra of compact distribution kernels and positive-scale contractions, gives
\begin{equation}
\boxed{
(Q+d_{\mathrm{CE}})I_{N,L}
 +\frac12\{I_{N,L},I_{N,L}\}_L
 +\hbar\Delta_LI_{N,L}=0.
}
\label{eq:hj-41}
\end{equation}
Its scalar component, using \eqref{eq:hj-37} and \eqref{eq:hj-39}, gives separately
\begin{equation}
\boxed{\Delta_LT_{n,N}^L=0\quad\text{for every }n.}
\label{eq:hj-42}
\end{equation}
Its quadratic components are
\begin{equation}
(Q+d_{\mathrm{CE},1})T_{n,N}^L
 +d_{\mathrm{CE},2}T_{n-1,N}^L
 +\frac12\sum_{i+j=n}
     \{T_{i,N}^L,T_{j,N}^L\}_L=0.
\label{eq:hj-43}
\end{equation}
Here $T_{0,N}^L=0$, and the bracket sum uses $i,j\geq1$.

Thus no central extension is needed. In this example, the classical tree components themselves are quantum current components.

\section{Quantum currents and every higher identity}

The master equation packages infinitely many multilinear identities.
Their target differential includes $\hbar\Delta_L$, even though the
current components themselves have no positive power of $\hbar$.
Fix the suspension signs to recover each identity from that equation.

\begin{theorem}\label{thm:hj-finite-current}
For every $N\geq1$ and $L>0$, the maps \eqref{eq:hj-46} form a
nontrivial quantum $L_\infty$ current morphism from $\mathfrak a_N$.
Its source has no central extension. Its values are compact quadratic
distribution kernels, with continuity on the stated fixed-support families.
Pullback along $q_N$ gives a morphism from the source
\eqref{eq:hj-27}, and from the projection of \eqref{eq:hj-29}.
\end{theorem}
\begin{proof}
Let
\[
\mathcal L_{N,L}
 =\mathcal O_{\mathrm{dist},N}[[\hbar]][-1],
\qquad |uF|=|F|+1,
\]
with
\begin{equation}
d_{\mathcal L}(uF)=-u(Q+\hbar\Delta_L)F,
\qquad
[uF,uG]_L=u\{F,G\}_L.
\label{eq:hj-44}
\end{equation}

For homogeneous \(\alpha_i\in\mathfrak a_N^{p_i}\), define
\begin{equation}
\begin{aligned}
K_{n,N,L}(\alpha_1,\ldots,\alpha_n)
={}&\int_{X^n}P_{0,L,12}\cdots P_{0,L,n-1,n}\\
&\quad\cdot
\beta_1\rho_N(\alpha_1)_1\cdots
\rho_N(\alpha_n)_n\gamma_n\,
\Omega_1\cdots\Omega_n.
\end{aligned}
\label{eq:hj-45}
\end{equation}
Let
\[
\operatorname{Alt}K(\alpha_1,\ldots,\alpha_n)
 =\sum_{\sigma\in S_n}
   \operatorname{sgn}(\sigma)\epsilon(\sigma;p)
   K(\alpha_{\sigma(1)},\ldots,\alpha_{\sigma(n)}).
\]
The maps are
\begin{equation}
\boxed{
\begin{aligned}
&F_{n,N,L}(\alpha_1,\ldots,\alpha_n)\\
&\qquad={}
(-1)^{p_1+\cdots+p_n+1}
u\,\operatorname{Alt}K_{n,N,L}(\alpha_1,\ldots,\alpha_n).
\end{aligned}
}
\label{eq:hj-46}
\end{equation}
They have degree \(1-n\). There is no omitted wheel term: its polarized value vanishes by \eqref{eq:hj-37}.
Indeed, if the two remaining matter forms have degrees
$q_\beta,q_\gamma$, the integral requires
\[
q_\beta+q_\gamma+\sum_i p_i+(n-1)=2n.
\]
The observable degree is $1-q_\beta-q_\gamma=\sum_i p_i-n$.
The shift by $u$ gives the asserted map degree.
Here $\epsilon(\sigma;p)$ is the Koszul sign for the input degrees
$p_i$, and the displayed antisymmetrization has no factorial divisor.

The pullback from formal Hamiltonian currents is explicitly
\begin{equation}
F_{n,N,L}^+
   =F_{n,N,L}\circ q_N^{\otimes n}.
\label{eq:hj-47}
\end{equation}

For a graded space $W$, use $S^c(W)=\bigoplus_{m\geq0}S^m(W)$.
Multiplication satisfies $xy=(-1)^{|x||y|}yx$.
The coaugmentation is $1$, the counit projects to length zero,
and the coproduct sums over all two-block unshuffles:
\[
\Delta(x_1\cdots x_n)=
\sum_{i=0}^n\sum_{\sigma\in\operatorname{Sh}(i,n-i)}
\epsilon_s(\sigma)
(x_{\sigma(1)}\cdots x_{\sigma(i)})
\otimes(x_{\sigma(i+1)}\cdots x_{\sigma(n)}).
\]
An unshuffle preserves the order inside each block, and
$\epsilon_s$ is the Koszul sign for the shifted degrees.
The two empty-block terms give the counit identities.
Splitting into three blocks proves coassociativity.

Set $x_i=s\alpha_i$, with $|x_i|=p_i-1$.
On $S^c(\mathfrak a_N[1])$, put
\begin{equation}
q_1(s\alpha)=s(d\alpha),\qquad
q_2(s\alpha\,s\eta)=(-1)^{|\alpha|}s[\alpha,\eta].
\label{eq:hj-48}
\end{equation}
Extend them by the signed unshuffle coderivation formulas. Jacobi and the differential derivation identity give \(q^2=0\).
Explicitly, a component $q_j$ contributes
\[
\sum_{\sigma\in\operatorname{Sh}(j,n-j)}\epsilon_s(\sigma)
q_j(x_{\sigma(1)}\cdots x_{\sigma(j)})
x_{\sigma(j+1)}\cdots x_{\sigma(n)}
\]
on a length-$n$ word, and $q(1)=0$ on the empty word.
The binary cancellation is the differential derivation identity;
the ternary cancellation is Jacobi. These are the only nonzero
corestrictions of $q^2$, so
Lemma~\ref{lem:aac-symmetric-coderivations} gives the full square-zero identity.
The convention in \eqref{eq:hj-48} uses $+sd$ as its unary component.

Identify \(s'(uF)=F\) on the target suspension. Its operations are
\begin{equation}
q'_1(F)=-(Q+\hbar\Delta_L)F,\qquad
q'_2(F,G)=(-1)^{|F|+1}\{F,G\}_L.
\label{eq:hj-49}
\end{equation}
The BV identities give \((q')^2=0\).

Set
\begin{equation}
f_{n,N,L}(s\alpha_1\cdots s\alpha_n)
 =
(-1)^{\sum_i(n-i)p_i}s'F_{n,N,L}(\alpha_1,\ldots,\alpha_n).
\label{eq:hj-50}
\end{equation}
These maps have degree zero and are graded symmetric.
The sign can be checked directly from an ordered background coefficient.
Put
\[
B_n(p)=\sum_i(n-i)p_i,\qquad
E_n(p)=\sum_{i<j}(p_i-1)(p_j-1),\qquad
\delta_n(p)=B_n(p)+E_n(p).
\]
Moving the formal parameters to the left of the ordered path gives
\[
M_n(p)=n\sum_i(1-p_i)+\sum_{i<j}p_i(1-p_j).
\]
The identity
\[
\frac{(n-1)(n-2)}2+M_n(p)+\delta_n(p)
\equiv\sum_i p_i+1\pmod2
\]
gives \eqref{eq:hj-46}. Interchanging adjacent degrees $p,q$
contributes $(-1)^{1+pq}$, which gives its graded antisymmetry.
The graded-dual product convention adds $(-1)^{E_n(p)}$ on evaluation,
so \eqref{eq:hj-50} is precisely
\[
f_n(x_1\cdots x_n)=(-1)^{E_n(p)}
[t_1\cdots t_n]T_{n,N}^L\left(\sum_i t_i\alpha_i\right).
\]

Their full coalgebra map is
\begin{equation}
\begin{aligned}
\mathbf f(x_1\cdots x_n)
={}&
\sum_{k=1}^n\frac1{k!}
\sum_{\substack{n_1+\cdots+n_k=n\\n_j\geq1}}
\sum_{\sigma\in\operatorname{Sh}(n_1,\ldots,n_k)}
\epsilon_s(\sigma)\\
&\qquad\cdot
f_{n_1}(x_{\sigma(1)},\ldots)\cdots
f_{n_k}(\ldots,x_{\sigma(n)}),
\end{aligned}
\label{eq:hj-51}
\end{equation}
with \(\mathbf f(1)=1\). This finite formula is well-defined on each input arity.
The successive arguments of the factor $f_{n_j}$ are exactly the
$j$th unshuffle block. The factor $1/k!$ removes the ordering of
the $k$ output blocks, with their graded symmetric signs.

Polarizing \eqref{eq:hj-41}, with the ghost-first rule \eqref{eq:hj-30}, gives
\begin{equation}
q'\mathbf f=\mathbf f q.
\label{eq:hj-52}
\end{equation}
More explicitly, the arity-\(n\) identity is
\begin{equation}
\begin{aligned}
q'_1f_n(x_1,\ldots,x_n)
&+\frac12\sum_{i+j=n}
 \sum_{\sigma\in\operatorname{Sh}(i,j)}
 \epsilon_s(\sigma)
 q'_2\bigl(f_i(x_{\sigma(1)},\ldots),
           f_j(\ldots,x_{\sigma(n)})\bigr)\\
&=\sum_{k=1}^n
 (-1)^{|x_1|+\cdots+|x_{k-1}|}
 f_n(x_1,\ldots,q_1x_k,\ldots,x_n)\\
&\quad+\sum_{\sigma\in\operatorname{Sh}(2,n-2)}
 \epsilon_s(\sigma)
 f_{n-1}\bigl(q_2(x_{\sigma(1)},x_{\sigma(2)}),
                   x_{\sigma(3)},\ldots,x_{\sigma(n)}\bigr).
\end{aligned}
\label{eq:hj-53}
\end{equation}
The binary target sum uses $i,j\geq1$.
The final source sum is absent when $n=1$.

For completeness, the passage from \eqref{eq:hj-41} to these signs
can be tested on a homogeneous suspended word $v$ of degree $r$.
Let $h_J(v)$ denote graded-dual evaluation of an interaction $J$,
including the product sign just specified. For $D_L=Q+\hbar\Delta_L$,
\[
h_{D_LJ}(v)=(-1)^rD_Lh_J(v),\qquad
h_{d_{\mathrm{CE}}J}(v)=(-1)^rh_J(qv).
\]
The first sign is the coefficient tensor differential sign.
The second follows from
$(d_{\mathrm{CE}}\phi)(v)=-(-1)^{|\phi|}\phi(qv)$.
For coefficients placed before observables, the bracket is
\[
\{\phi F,\psi G\}_L
=(-1)^{(|F|+1)|\psi|}\phi\psi\{F,G\}_L.
\]
For a split of degrees $r_1,r_2$, the extension sign and the
dual-product sign multiply to $(-1)^{r_2}$.
Consequently, if
$\mathcal R=(D_L+d_{\mathrm{CE}})J+\tfrac12\{J,J\}_L$, then
\[
(q'\mathbf f-\mathbf f q)_{\mathrm{corestr}}(v)
 =(-1)^{r+1}h_{\mathcal R}(v).
\]
Equation~\eqref{eq:hj-41} sets the right side to zero.

Expanding \eqref{eq:hj-51} by partitions proves its coalgebra identity.
Both $q'\mathbf f$ and $\mathbf f q$ are coderivations along
$\mathbf f$. Their difference is determined by its projection to
the cogenerators: induction on input arity reduces the difference
to a primitive element. In characteristic zero, the only primitives
of $S^c(W)$ are $W$. Indeed, multiplication composed with the component
\[
S^n(W)\longrightarrow W\otimes S^{n-1}(W)
\]
of the reduced coproduct is \(n\) times the identity. Thus vanishing of the cogenerator projection proves \eqref{eq:hj-52} in every arity.
For $n\geq2$ this kills a primitive's length-$n$ component.
A primitive scalar $\lambda1$ must satisfy $\lambda=2\lambda$
and therefore vanishes as well. The induction begins at the empty
word, which both coderivations kill.

Equations \eqref{eq:hj-46}--\eqref{eq:hj-53} give the actual quantum \(L_\infty\) morphism from each finite source, and its pullback from \eqref{eq:hj-27} or \eqref{eq:hj-29}.
Nontriviality is already visible in the $H_1$ block at every $N$:
$\rho_N(e)y=x$. Choose a compact $(0,2)$ source form whose integral
against $\Omega$ is nonzero, and smooth matter coefficients selecting
$\beta_x$ and $\gamma_y$. Its linear current is nonzero.
The fixed-support continuity follows from \eqref{eq:hj-34} and its
finite polarization. This completes the construction.
\end{proof}

\section{What a matter projection fails to preserve}

The quotient modules form a compatible representation system.
A current theory also pairs a matter module with its dual.
The variance of that dual, and the survival of scalar contractions,
must therefore be checked before taking a quantum limit.
The maps
\begin{equation}
\pi_{N+1,N}:\mathfrak g_{N+1}\to\mathfrak g_N,
\qquad
p_{N+1,N}:V_{N+1}\to V_N
\label{eq:hj-54}
\end{equation}
are the quotient projections. They satisfy
\begin{equation}
p_{N+1,N}\rho_{N+1}(f)
 =\rho_N(\pi_{N+1,N}f)\,p_{N+1,N}.
\label{eq:hj-55}
\end{equation}
This follows directly by taking the same Poisson bracket before either quotient.

The module projection has no \(H_+\)-equivariant section. Suppose \(s\) were such a section. Then
\[
s(y)=y+w,\qquad w\in H_{N+1}.
\]
Take \(f=x^{N+2}\). It acts trivially on \(V_N\), but
\[
\rho_{N+1}(f)s(y)
 =(N+2)x^{N+1}\ne0.
\]
The bracket with \(w\) has degree \(2N+1\), hence vanishes modulo \(\mathfrak m^{N+2}\). This contradicts equivariance.

The canonical dual map is
\begin{equation}
p_{N+1,N}^\vee:V_N^\vee\hookrightarrow V_{N+1}^\vee,
\label{eq:hj-56}
\end{equation}
in the opposite direction. There is no automatic paired matter projection supplied by \eqref{eq:hj-55}.

Using the homogeneous splitting, one can nevertheless delete every \(\beta\)- and \(\gamma\)-coordinate of internal degree \(N+1\). Call this observable restriction \(r_N\). It commutes with \(Q\), but it does not preserve BV contractions.

Choose a removed pair of internal dual basis vectors. Choose compact linear observables \(b,g\) in paired form components for which the positive-scale heat contraction is nonzero. Such choices exist because the heat kernel and internal evaluation pairing are nonzero. Normalize them so that
\[
\{b,g\}_L=1.
\]
Then
\[
r_N(b)=r_N(g)=0,\qquad r_N(1)=1,
\]
and therefore
\begin{equation}
r_N\{b,g\}_L=1
 \ne \{r_Nb,r_Ng\}_L.
\label{eq:hj-57}
\end{equation}
Thus deletion is not a morphism of the ambient BV algebras or their shifted Lie algebras.

The opposite inclusion of lower modes preserves the free contractions, but does not intertwine the current maps. Already for \(N=1\to2\),
\begin{equation}
\rho_1(x^3)=0,\qquad
J_2(\alpha x^3)
 =3\int_X\beta_{x^2}\,\alpha\,\gamma_y\,\Omega,
\label{eq:hj-58}
\end{equation}
which is generally nonzero.

These are concrete failures of the canonical ambient transition maps. They do not rule out every possible construction involving additional fields, homotopies, or integration over discarded modes.

\section{An inverse system of current Lie algebras}

The currents retain a property absent from the discarded linear pair:
their mixed quadratic kernels never lower internal degree.
This bounds every intermediate degree in a single-contraction bracket,
and permits a smaller inverse system despite \eqref{eq:hj-57}.

At fixed \(N,L\), consider mixed quadratic observables
\[
F=\sum_{a,b}\beta_a\,K_{ab}\,\gamma_b,
\]
where \(K_{ab}\) denotes a compact spatial distribution kernel, including its form and density labels. Call such a kernel degree-nondecreasing when
\begin{equation}
K_{ab}=0\quad\text{if }d(a)<d(b).
\label{eq:hj-59}
\end{equation}
Here $a$ is an output monomial label, $b$ an input monomial label,
and $d(a),d(b)$ are their positive total degrees in the internal variables.
Every tree kernel \eqref{eq:hj-45} has this property, by \eqref{eq:hj-6}.

The bracket of two mixed quadratic observables is again mixed quadratic. It contracts one \(\beta\)-slot with one \(\gamma\)-slot. Its internal matrix terms are compositions, with the smooth heat kernel inserted between their spatial slots.

For a matrix coefficient with low input degree \(b\) and low output degree \(a\), a nonzero intermediate degree \(c\) must satisfy
\begin{equation}
d(b)\leq d(c)\leq d(a).
\label{eq:hj-60}
\end{equation}
Hence, when \(d(a),d(b)\leq N\), every intermediate degree is already at most \(N\). Deleting higher modes therefore preserves the bracket on degree-nondecreasing mixed quadratics:
\begin{equation}
r_N\{F,G\}_L=\{r_NF,r_NG\}_L.
\label{eq:hj-61}
\end{equation}
The graded signs are unchanged, since deletion does not reorder any surviving slot. The same argument proves closure of condition \eqref{eq:hj-59} under the bracket. Spatial differentiation also preserves it.

Define \(S_{N,L}^0\) to be the smallest graded Lie subspace of
\(\mathcal O_{\mathrm{dist},N}[-1]\) which:

\begin{enumerate}
\item contains every value of every map \(F_{n,N,L}\);
\item is preserved by \(uF\mapsto-uQF\).
\end{enumerate}

Every element is a degree-nondecreasing mixed quadratic. Equation \eqref{eq:hj-42} kills the BV Laplacian on the generators. The identities
\[
Q\Delta_L+\Delta_LQ=0
\]
and the BV derivation identity for \(\Delta_L\) acting on the odd bracket then give
\begin{equation}
\Delta_LF=0\qquad(uF\in S_{N,L}^0).
\label{eq:hj-62}
\end{equation}
Therefore
\begin{equation}
S_{N,L}=S_{N,L}^0[[\hbar]]
\label{eq:hj-63}
\end{equation}
is a differential graded Lie subalgebra of the quantum target \eqref{eq:hj-44}. Its differential is simply \(-uQ\), because its \(\Delta_L\) term vanishes.

\begin{theorem}\label{thm:hj-current-limit}
Restriction to internal degrees at most $N$ gives surjective
differential graded Lie maps $S_{N+1,L}\to S_{N,L}$.
The current components are compatible with these maps and the
source quotient maps. They define the $L_\infty$ morphism
\eqref{eq:hj-68} into the algebraic inverse limit.
\end{theorem}
\begin{proof}
For compatible source inputs, deleting higher modes in \eqref{eq:hj-45} gives exactly
\begin{equation}
r_NF_{n,N+1,L}(\alpha_1,\ldots,\alpha_n)
 =
F_{n,N,L}(\pi_{N+1,N}\alpha_1,\ldots,
                        \pi_{N+1,N}\alpha_n).
\label{eq:hj-64}
\end{equation}
Indeed, low matrix coefficients of the product use only low intermediate degrees by \eqref{eq:hj-60}. The scalar propagators and integration domains are identical.

Equations \eqref{eq:hj-61}, \eqref{eq:hj-62}, and commutation with \(Q\) now prove that
\begin{equation}
\boxed{
r_N:S_{N+1,L}\longrightarrow S_{N,L}
}
\label{eq:hj-65}
\end{equation}
is a differential graded Lie map.

It is surjective. The source projection \(\mathfrak a_{N+1}\to\mathfrak a_N\) is surjective: lift each finite-dimensional coefficient by its homogeneous polynomial representative. Thus every generating current at level \(N\) lifts through \eqref{eq:hj-64}. Lift a finite expression in generators, their differentials, and their brackets by lifting its generators. Formal \(\hbar\)-series lift coefficientwise.

Consequently,
\begin{equation}
S_{\infty,L}=\varprojlim_NS_{N,L}
\label{eq:hj-66}
\end{equation}
is a well-defined differential graded Lie algebra. Its operations are coordinatewise. No additional completeness of the generated distribution subspaces is asserted.
This is an algebraic inverse limit, with formal $\hbar$ coefficients
treated coefficientwise. Its projective topology is specified by the
finite-level distribution coordinates. The analytic estimates remain
the fixed-compact-support estimates at each finite level.
They do not assert unrestricted global joint continuity, completeness
of $S_{N,L}^0$, or one compact kernel on an infinite-dimensional
matter space.

For \(\alpha_i=(\alpha_{i,N})_N\in\mathfrak a_+\), define
\begin{equation}
\boxed{
F_{n,\infty,L}(\alpha_1,\ldots,\alpha_n)
 =
\bigl(F_{n,N,L}(\alpha_{1,N},\ldots,\alpha_{n,N})\bigr)_N.
}
\label{eq:hj-67}
\end{equation}
Compatibility follows from \eqref{eq:hj-64}. Every all-arity identity holds after every coordinate projection, hence holds in the inverse limit. Thus \eqref{eq:hj-67} is an actual \(L_\infty\) morphism
\begin{equation}
\mathfrak a_+\longrightarrow S_{\infty,L}.
\label{eq:hj-68}
\end{equation}
\end{proof}

This inverse system preserves the differential, the single-contraction BV
bracket, the vanishing BV Laplacian on the current Lie space, and every
current component. Products can leave that Lie space. To retain them,
we next form the algebra generated by its actual distribution kernels
and check the BV identities there. The full ambient counterexample
\eqref{eq:hj-57} remains a separate restriction.

\section{Products of currents and the BV operator}
\label{sec:hj-bv-products}

A current map gives individual insertions and their higher relations.
Products of insertions require an algebra as well. The Lie space
$S_{N,L}^0$ consists of mixed quadratics, so it cannot contain a
general product of two of its elements. Its vanishing Laplacian
does not justify setting the Laplacian to zero on those products.
A contraction can join two distinct current factors.

Fix $L>0$ throughout the following construction. All degrees are
cohomological degrees of unshifted observables, and $\hbar$ has
degree zero. The BV convention is
\begin{equation}\label{eq:hj-bv-bracket}
\{F,G\}_L=(-1)^{|F|}
 \left(\Delta_L(FG)-(\Delta_LF)G-(-1)^{|F|}F\Delta_LG\right).
\end{equation}
Thus $\Delta_L$ and the bracket have degree one. Products obey
$FG=(-1)^{|F||G|}GF$. The algebra unit has degree zero.

First suppose $F,G,H$ are homogeneous and individually
$\Delta_L$-closed. Write $f=|F|$, $g=|G|$, $h=|H|$.
The binary product identity gives
\begin{equation}\label{eq:hj-bv-pair}
\boxed{\Delta_L(FG)=(-1)^f\{F,G\}_L.}
\end{equation}
The first-slot product rule for the odd bracket is
\[
\{FG,H\}_L
 =F\{G,H\}_L+(-1)^{g(h+1)}\{F,H\}_LG.
\]
Applying \eqref{eq:hj-bv-bracket} to $(FG)H$ therefore gives
\begin{equation}\label{eq:hj-bv-triple}
\boxed{\begin{aligned}
\Delta_L(FGH)
={}&(-1)^f\{F,G\}_LH
   +(-1)^{f+g}F\{G,H\}_L\\
 &+(-1)^{f+gh}\{F,H\}_LG.
\end{aligned}}
\end{equation}
The order of the remaining factors is part of this formula.
If the three pair brackets are also $\Delta_L$-closed, a second
application gives
\begin{equation}\label{eq:hj-bv-triple-square}
\begin{aligned}
\Delta_L^2(FGH)
={}&(-1)^{g+1}\{\{F,G\}_L,H\}_L
   +(-1)^g\{F,\{G,H\}_L\}_L\\
 &+(-1)^{gh+h+1}\{\{F,H\}_L,G\}_L=0.
\end{aligned}
\end{equation}
The final equality is precisely the odd Jacobi identity, with
the displayed factor orders. It verifies the ternary signs
without assuming that products remain in the kernel of $\Delta_L$.

The first internal jet supplies a finite sign calculation.
Take odd variables $\beta_x,\beta_y$ of degree one and even
variables $\gamma_x,\gamma_y$ of degree minus two. Put
\[
\Delta=-\partial_{\beta_x}\partial_{\gamma_x}
       -\partial_{\beta_y}\partial_{\gamma_y},
\qquad
F=\beta_x\gamma_y,\qquad G=\beta_y\gamma_x.
\]
The paired degrees sum to minus one, so $\Delta$ has degree one.
Left differentiation in the odd variables gives
\begin{equation}\label{eq:hj-bv-coordinate}
\Delta F=\Delta G=0,\qquad
\Delta(FG)=\beta_x\gamma_x-\beta_y\gamma_y\ne0.
\end{equation}
The two diagonal terms have opposite Laplacians, so the last
expression is $\Delta$-closed. These are the internal matrices
$e,f,h$ from the first jet. A spatial current calculation below
will exhibit the same nonzero product contraction in the actual
distribution algebra.

\section{The algebra generated by actual current kernels}
\label{sec:hj-bv-algebra}

To retain multiplication, use the current kernels as elements of
their existing observable algebra. This preserves relations that
would disappear if the currents were replaced by independent symbols.
Define the unshifted current space
\begin{equation}\label{eq:hj-bv-current-space}
C_{N,L}=\{F\in\mathcal O_{\mathrm{dist},N}:uF\in S_{N,L}^0\}.
\end{equation}
The preceding current construction proves
\begin{equation}\label{eq:hj-bv-current-closure}
QC_{N,L}\subset C_{N,L},\qquad
\{C_{N,L},C_{N,L}\}_L\subset C_{N,L},\qquad
\Delta_LC_{N,L}=0.
\end{equation}
Every element is a mixed quadratic distribution kernel whose
internal matrix coefficients satisfy \eqref{eq:hj-59}.

Let $\operatorname{Sym}_{\mathbb C}(C_{N,L})$ be the algebraic
graded symmetric algebra, with its unit in symmetric power zero.
Multiplication of actual observables gives an algebra map
\[
\mu_N:\operatorname{Sym}_{\mathbb C}(C_{N,L})
          \longrightarrow\mathcal O_{\mathrm{dist},N}.
\]
Define
\begin{equation}\label{eq:hj-bv-image}
\boxed{A^0_{N,L}=\operatorname{im}\mu_N.}
\end{equation}
An element is a finite sum of finite products of elements of
$C_{N,L}$. The empty product maps to the scalar observable $1$.
The products use external tensor products of distributions,
followed by the graded symmetric coinvariants of their matter slots.
No product of distributions is restricted to a spatial diagonal.

Write $\mathcal I_{N,L}=\ker\mu_N$. Then
\begin{equation}\label{eq:hj-bv-relations}
A^0_{N,L}\cong
\operatorname{Sym}_{\mathbb C}(C_{N,L})/\mathcal I_{N,L}.
\end{equation}
Injectivity of $\mu_N$ is neither needed nor asserted.
All polynomial relations already present among the distribution
observables remain in this quotient.

There is an additional grading by the number of current factors:
\begin{equation}\label{eq:hj-bv-factor-grading}
A^0_{N,L}=\bigoplus_{r\geq0}A^{[r]}_{N,L},\qquad
A^{[r]}_{N,L}=\mu_N\bigl(\operatorname{Sym}^r C_{N,L}\bigr).
\end{equation}
This is not the cohomological grading. Each current factor has
two matter slots, so $A^{[r]}_{N,L}$ lies in observable arity $2r$.
Different arities are distinct direct summands of
$\mathcal O_{\mathrm{dist},N}$. Thus the sum in
\eqref{eq:hj-bv-factor-grading} is direct, and the relation ideal
$\mathcal I_{N,L}$ is homogeneous in this factor grading.
Indeed, if a finite sum of homogeneous components maps to zero,
its components map into distinct observable arities and hence each
maps to zero. Thus every homogeneous component of a relation is
again a relation.

A unital differential BV algebra here is a unital graded commutative
algebra with a degree-one derivation $Q$ and a degree-one operator
$\Delta$ of order at most two, satisfying
$Q^2=\Delta^2=Q\Delta+\Delta Q=0$ and $Q1=\Delta1=0$.
Its bracket is derived by \eqref{eq:hj-bv-bracket}.

\begin{theorem}\label{thm:hj-bv-finite}
The actual image $A^0_{N,L}$ is a unital differential BV subalgebra
of $\mathcal O_{\mathrm{dist},N}$. Its operators are the restrictions
of $Q$ and $\Delta_L$. It is the smallest such subalgebra containing
every unshifted current component.
\end{theorem}
\begin{proof}
The ambient operator $Q$ is a derivation and preserves $C_{N,L}$.
Hence $Q(A^{[r]}_{N,L})\subset A^{[r]}_{N,L}$.
The second-slot product rule is
\[
\{F,GH\}_L=\{F,G\}_LH
 +(-1)^{(|F|+1)|G|}G\{F,H\}_L.
\]
Together with the first-slot rule and
\eqref{eq:hj-bv-current-closure}, it proves
\begin{equation}\label{eq:hj-bv-bracket-grading}
\{A^{[r]}_{N,L},A^{[s]}_{N,L}\}_L
          \subset A^{[r+s-1]}_{N,L}\qquad(r,s\geq1).
\end{equation}
Brackets with the scalar summand vanish.

For homogeneous $F_1,\ldots,F_r\in C_{N,L}$, put
\begin{equation}\label{eq:hj-bv-pair-sign}
\epsilon_{ij}=(-1)^{
 |F_i|\sum_{k<i}|F_k|
 +|F_j|\sum_{\substack{k<j\\k\ne i}}|F_k|}.
\end{equation}
This is the Koszul sign moving $F_i,F_j$, in that order, to the
front. The full product formula is
\begin{equation}\label{eq:hj-bv-all-products}
\boxed{\begin{aligned}
\Delta_L(F_1\cdots F_r)
={}&\sum_{i<j}(-1)^{|F_i|}\epsilon_{ij}\{F_i,F_j\}_L\\
&\hspace{9mm}\cdot
F_1\cdots\widehat F_i\cdots\widehat F_j\cdots F_r.
\end{aligned}}
\end{equation}
The bracket stands first. The other factors retain their original
order, and a hat denotes omission.

For two factors this is \eqref{eq:hj-bv-pair}.
Induct on $r$ by applying \eqref{eq:hj-bv-bracket} to
$F_1(F_2\cdots F_r)$. Terms from the Laplacian of the second
factor account for pairs not containing $1$.
The bracket biderivation accounts for the pairs $(1,j)$.
Moving each chosen pair to the front gives
$\epsilon_{ij}$, and contracting that ordered pair gives
$(-1)^{|F_i|}\{F_i,F_j\}_L$.
These exhaust the terms; single-factor terms vanish by
\eqref{eq:hj-bv-current-closure}. This proves the formula.
In particular,
\begin{equation}\label{eq:hj-bv-lowering}
\Delta_L(A^{[r]}_{N,L})\subset A^{[r-1]}_{N,L},\qquad
\Delta_L(A^{[0]}_{N,L}\oplus A^{[1]}_{N,L})=0.
\end{equation}

These operations are restrictions of the actual ambient operators.
They therefore satisfy
\[
Q^2=\Delta_L^2=Q\Delta_L+\Delta_LQ=0,\qquad
Q1=\Delta_L1=0,
\]
and $\Delta_L$ has differential-operator order at most two.
This order is its order relative to multiplication, not its
cohomological degree or its change of factor grading.

Relations cause no further obstruction. Applying an ambient
operator to a zero distribution gives zero. Equivalently, the
formulas for $Q$, the bracket, and $\Delta_L$ on the symmetric
algebra preserve $\mathcal I_{N,L}$ and descend to its actual image.
Thus $A^0_{N,L}$ is a unital differential BV algebra.

Any unital differential BV subalgebra containing the current
components contains their $Q$-images and derived brackets.
It therefore contains $C_{N,L}$ and all finite products defining
$A^0_{N,L}$. This proves the asserted minimality within the fixed
ambient distribution algebra. No maximality or topological closure
is asserted.
\end{proof}

\section{Why restriction preserves these products}
\label{sec:hj-bv-restriction}

The ambient restriction $r_N$ already preserves multiplication:
deletion acts independently on each matter slot and commutes with
external tensor product and graded permutations. Thus
\begin{equation}\label{eq:hj-bv-algebra-restriction}
r_N(FG)=(r_NF)(r_NG),\qquad r_N1=1.
\end{equation}
Its failure in \eqref{eq:hj-57} concerns contractions.
Two specific facts about the current space remove that failure
from its generated algebra.

First, in a cross contraction with surviving output $a$, input $b$,
and internal label $c$, a coefficient has the form
$K^F_{ac}H_LK^G_{cb}$, or its reverse orientation. The two
nondecreasing matrices impose
\[
d(b)\leq d(c)\leq d(a).
\]
Low surviving labels therefore force a low contracted label.
This is the bracket compatibility already proved in
\eqref{eq:hj-61}. Second, every current-space element has zero
self-contraction by \eqref{eq:hj-bv-current-closure}.
Nondecreasing degree alone would not imply this second fact.

\begin{proposition}\label{prop:hj-bv-restriction}
The restrictions induce surjective morphisms of unital differential
BV algebras
\[
r_N:A^0_{N+1,L}\longrightarrow A^0_{N,L}.
\]
\end{proposition}
\begin{proof}
The current restriction is a surjection
$C_{N+1,L}\to C_{N,L}$ preserving $Q$ and the bracket, by
Theorem~\ref{thm:hj-current-limit}.
Equation~\eqref{eq:hj-bv-algebra-restriction} therefore maps the
upper generated algebra into the lower one.
Lift a finite polynomial by choosing lifts of its finitely many
current-space factors. This proves surjectivity.

Apply \eqref{eq:hj-bv-all-products} to a product of upper
current-space elements. Every pair bracket is preserved by
$r_N$, and restriction introduces no permutation of the surviving
factors. Consequently,
\begin{equation}\label{eq:hj-bv-delta-restriction}
\boxed{r_N\Delta_L=\Delta_Lr_N\quad\hbox{on }A^0_{N+1,L}.}
\end{equation}
Commutation with $Q$ follows from its ambient restriction identity.
This argument respects actual relations: an upper expression which
is zero as a distribution has zero ambient restriction.
It is not necessary to choose polynomial representatives consistently.
\end{proof}

Set
\begin{equation}\label{eq:hj-bv-hbar}
A_{N,L}=A^0_{N,L}[[\hbar]],\qquad |\hbar|=0,
\end{equation}
with the formal series taken degreewise in the cohomological grading.
For each fixed cohomological degree, coefficients lie in that
degree of $A^0_{N,L}$. The graded algebra has finite sums of
homogeneous cohomological components; no completion in that grading
is imposed. All operations extend coefficientwise, using finite
Cauchy sums for products. A formal series lifts by lifting each
coefficient, so
\begin{equation}\label{eq:hj-bv-formal-restriction}
\boxed{r_N:A_{N+1,L}\twoheadrightarrow A_{N,L}}
\end{equation}
is a surjective morphism of unital differential BV algebras over
$\mathbb C[[\hbar]]$.

\section{A nonzero contraction of actual current products}
\label{sec:hj-bv-nonzero-product}

The restricted algebra has a nontrivial quantum operator, despite
the vanishing of every complete scalar wheel. Put
$D_L=Q+\hbar\Delta_L$. Its square is zero. For homogeneous
$F,G\in A_{N,L}$, its product defect is
\begin{equation}\label{eq:hj-bv-quantum-product}
\begin{aligned}
D_L(FG)-(D_LF)G-(-1)^{|F|}FD_LG
   =\hbar(-1)^{|F|}\{F,G\}_L.
\end{aligned}
\end{equation}
This follows from the derivation rule for $Q$ and
\eqref{eq:hj-bv-bracket}. Thus $D_L$ is not in general a
derivation of multiplication.

To exhibit a nonzero defect in the actual current algebra, take
$N=1$. Choose disjoint open balls $U,V\subset X$ and nonzero,
nonnegative compact smooth functions $\chi_U,\chi_V$ supported
in them. Use the form-degree-zero sources
\[
\alpha=\chi_Ue,\qquad \eta=\chi_Vf,
\qquad F=J(\alpha),\qquad G=J(\eta).
\]
Both currents belong to $C_{1,L}$, have cohomological degree
minus one, and are $\Delta_L$-closed.
The fixed-scale contraction formula is
\begin{equation}\label{eq:hj-bv-current-bracket}
\{F,G\}_L
 =-K_{H_L}(\alpha,\eta)+K_{H_L}(\eta,\alpha),
\end{equation}
where the kernel in \eqref{eq:hj-45} at $n=2$ uses $H_L$ in
place of $P_{0,L}$. This is the positive-scale current bracket
derived in Chapter~\ref{chap:quantum-current-anomaly}.

Select only the external components
\[
\beta_x(z)=b\,\theta_U(z),\qquad
\gamma_x(w)=g\,\theta_V(w)\,
                     d\bar w_1\wedge d\bar w_2,
\]
where $|b|=1$, $|g|=-2$. Choose nonnegative compact
$\theta_U,\theta_V$ with
$\chi_U\theta_U\not\equiv0$ and $\chi_V\theta_V\not\equiv0$,
supported in $U,V$, respectively. All other matter components
are zero. The reversed kernel vanishes because its $\beta$-slot
lies in $V$. The first kernel survives because $(ef)_{xx}=1$.

The surviving heat component is
$d\bar z_1d\bar z_2d\bar w_1d\bar w_2$.
Each orientation identity contributes a factor $4$, so the
coefficient of $bg$ in that first kernel is
\begin{equation}\label{eq:hj-bv-positive-integral}
\frac1{4\pi^2L^2}\int_{X^2}
e^{-|z-w|^2/(4L)}
\chi_U(z)\theta_U(z)\chi_V(w)\theta_V(w)
\,dV_4(z)dV_4(w)>0.
\end{equation}
The integrand is nonnegative, and the two overlap conditions make
it positive on a set of positive measure. The constant is
$4\cdot4/(64\pi^2L^2)$, with $\Omega_z,\Omega_w$ on the right.
Since $|F|=-1$, \eqref{eq:hj-bv-pair} and
\eqref{eq:hj-bv-current-bracket} give
\begin{equation}\label{eq:hj-bv-nonzero-product}
\boxed{\Delta_L(FG)
 =K_{H_L}(\alpha,\eta)-K_{H_L}(\eta,\alpha)\ne0.}
\end{equation}
Thus the algebra generated by the actual currents has a nonzero
BV operator. Spatial separation does not make multiplication a
quantum chain map at a positive heat scale.

\section{The algebraic BV limit and its distribution topology}
\label{sec:hj-bv-limit}

The restriction identities allow all algebra operations to pass
to compatible finite-level coefficients. Define, degreewise,
\begin{equation}\label{eq:hj-bv-limit}
\boxed{A_{\infty,L}=\varprojlim_N A_{N,L}.}
\end{equation}
In a fixed cohomological degree, an element is a compatible array
\[
a=(a_{N,k})_{N\geq1,\ k\geq0},\qquad
a_{N,k}\in A^0_{N,L},\qquad r_Na_{N+1,k}=a_{N,k}.
\]
Each coefficient is a finite polynomial in current-space elements.
There is no uniform polynomial-degree or compact-support bound
across $N$ or $k$. In particular, no completion in observable
arity is taken at a fixed coefficient.

The operations are explicitly
\begin{equation}\label{eq:hj-bv-limit-operations}
\begin{split}
(ab)_{N,k}&=\sum_{i+j=k}a_{N,i}b_{N,j},\\
(Qa)_{N,k}&=Qa_{N,k},\qquad
(\Delta_La)_{N,k}=\Delta_La_{N,k},\\
(D_La)_{N,k}&=Qa_{N,k}+\Delta_La_{N,k-1},\qquad a_{N,-1}=0.
\end{split}
\end{equation}
The unit is the array of unit constant terms and zero positive
$\hbar$ coefficients. Every expression at a specified coordinate
is finite. All BV identities hold after each coordinate
projection, and those projections jointly distinguish elements.
This constructs a unital differential BV algebra, without
identifying it with observables on a single infinite-dimensional
matter space.

A precise distribution topology can be retained with this
algebraic construction. Give each observable-arity summand of
$\mathcal O_{\mathrm{dist},N}$ its strong distribution topology,
and give the algebraic sum over arities the locally convex
direct-sum topology. Give $A^0_{N,L}$ the induced subspace topology.
Then give $A_{N,L}$ the product topology in its $\hbar$ coefficients,
degreewise, and $A_{\infty,L}$ the subspace topology in
$\prod_NA_{N,L}$.

Explicit defining seminorms on the inverse limit have the form
\begin{equation}\label{eq:hj-bv-seminorms}
a\longmapsto
\sum_{m\geq0}\sup_{\phi\in B_m}
\left|\left\langle(a_{N,k})_m,\phi\right\rangle\right|.
\end{equation}
Here $N,k$ are fixed and each $B_m$ is a bounded family of smooth
tests in the corresponding graded symmetric test space.
The sum is finite on each $a_{N,k}$ because that coefficient has
finite observable arity. Finite sums of these seminorms give
the projective topology.

The maps $r_N,Q,\Delta_L$ are continuous linear maps for these
topologies. Differentiation, finite projections of internal labels,
and insertion of the fixed smooth heat kernel send bounded test
families to bounded test families. Transposition therefore gives
continuity on each strong distribution space. Aritywise continuity
extends to the locally convex direct sum, then to coefficient
products and the inverse limit.

For multiplication, if distributions $T,S$ have fixed compact
supports $K,K'$ and orders $p,q$, their external product obeys
\[
|(T\boxtimes S)(\phi)|\leq C_TC_S
\max_{\substack{|\mu|\leq p\\|\nu|\leq q}}
\sup_{K\times K'}|\partial_z^\mu\partial_w^\nu\phi|.
\]
Inserting a fixed smooth heat kernel gives the corresponding
finite-order bounds for brackets. These are fixed-support,
fixed-order estimates. No unrestricted joint continuity or
completeness of $A^0_{N,L}$ is asserted.

\section{Current maps into the BV algebra}
\label{sec:hj-bv-current-inclusion}

The algebra construction changes the target available for products,
but leaves every current component unchanged. The inclusion
\[
S_{N,L}\hookrightarrow A_{N,L}[-1],\qquad uF\longmapsto uF,
\]
is a differential graded Lie map. Its target differential is
$-uD_L$ and its bracket is $u\{F,G\}_L$.
On current-space elements $\Delta_LF=0$, so the differential
agrees with the existing $-uQ$.
The compatibility identity \eqref{eq:hj-64} therefore gives
\begin{equation}\label{eq:hj-bv-current-map}
\boxed{\mathfrak a_+\longrightarrow
S_{\infty,L}\hookrightarrow A_{\infty,L}[-1].}
\end{equation}
It uses exactly the components \eqref{eq:hj-67}.
Every higher identity holds coordinatewise as before.

The ambient counterexample remains: for a removed normalized
pair $b,g$, one has $r_Nb=r_Ng=0$ but
$r_N\{b,g\}_L=1$. In particular, the diagonal quadratic $bg$
is nondecreasing in internal degree, yet its Laplacian is not
preserved. The two conditions used here are bracket compatibility
on the current space and vanishing self-contractions there.
Neither condition should be dropped.

The resulting BV operations are the prescribed signed sums of
contractions. Their compatibility does not assert compatibility
of each individual graph contraction, ordinary matrix traces,
singular propagator operations on arbitrary distributions, or
Gaussian integration over discarded modes. Nor does the algebraic
inverse limit construct factorization maps or an analytic completion.
The trace and geometric comparisons can now be assessed separately
from this established algebra of current products.


\section{Compatible currents and incompatible traces}

The following data are exactly compatible as \(N\) increases:

\begin{itemize}
\item the Lie quotients \(\mathfrak g_N\);
\item the module projections \(V_{N+1}\to V_N\);
\item every fixed low matrix coefficient of an action word;
\item the tree current components under deletion;
\item the zero symmetric cubic tensor and zero complete wheel corrections;
\item the current Lie subalgebras and morphism \eqref{eq:hj-68};
\item the actual generated BV algebras and current map \eqref{eq:hj-bv-current-map}.
\end{itemize}

The ordinary finite-dimensional traces are not compatible. For instance,
\[
\operatorname{Tr}_{V_{N+1}}I
 -\operatorname{Tr}_{V_N}I=N+2,
\]
and
\begin{equation}
\operatorname{Tr}_{V_{N+1}}h^2
 -\operatorname{Tr}_{V_N}h^2
 =\frac{(N+1)(N+2)(N+3)}3.
\label{eq:hj-69}
\end{equation}
Equations \eqref{eq:hj-19} and \eqref{eq:hj-21} grow respectively as \(N^4/12\) and \(N^6/30\).

Therefore no scalar trace on the inverse system can both assign a single value to a compatible operator and equal its ordinary trace at every finite stage. The identity alone proves this obstruction.

The current BV limit preserves its prescribed algebra operations.
Complete scalar wheels vanish, while pair contractions can act
nontrivially on products, as \eqref{eq:hj-bv-nonzero-product} shows.
It does not identify all ordinary finite-dimensional matrix traces.

Likewise, the growth of \(\operatorname{Tr}(h^4)\) does not prove divergence of the complete four-wheel: that wheel is zero by \eqref{eq:hj-37}. A nonzero ordered color coefficient is not a complete physical graph sum. No physical quantum limit is inferred from the moment calculation alone.

\section{Translations and geometric jet prolongation}

The current inverse limit preserves origin-fixing internal symmetries.
Two separate comparisons remain: adding internal translations, and
allowing a Hamiltonian to move the spatial insertion point.
The full translation-containing algebra
\[
H=\mathbb C[[x,y]]/\mathbb C
\]
does not act on \(V_N\) by the quotient formula \eqref{eq:hj-2}. Indeed,
\[
\{x,y^{N+1}\}=(N+1)y^N
\]
shows that its translations do not preserve \(\mathfrak m^{N+1}\).

Deleting the linear Hamiltonian terms before applying \(\rho_N\) also fails to define a Lie algebra action. This projection is not a Lie algebra map:
\[
\{x,y^3\}=3y^2
\]
has a nonzero quadratic projection, while the projected first input is zero.

More generally, every continuous finite-dimensional representation of \(H\) is zero. Its kernel is an open ideal. An open ideal containing \(\mathfrak m^M\) contains
\[
I_y^Mf
 =\sum_{a+b\geq1}\frac{b!}{(b+M)!}
     f_{ab}x^ay^{b+M},
\]
and repeated brackets with \(x\) give
\[
(\operatorname{ad}_x)^MI_y^Mf=f.
\]
Thus the ideal is all of \(H\). If the finite-dimensional target uses its usual complex topology, continuity still yields an open kernel: an open linear subspace mapped into a bounded ball must map to zero.

The nontrivial jet representations therefore have no continuous extension to the full \(H\). The full operator trace obstructions remain separate and unchanged.

There is also a geometric distinction. The preceding quantum theory treats \(x,y\) as internal formal variables and \(z_1,z_2\) as spatial coordinates. Its matrix symmetry acts on matter without differentiating its spatial argument.

An explicit comparison with geometric Hamiltonian action requires jet prolongation. For actual holomorphic functions \(F,g\) on a spatial domain \(U\subset\mathbb C^2_z\), introduce internal jet variables \(\xi_1,\xi_2\) and put
\[
j_N^+g(z,\xi)
 =\sum_{1\leq|\alpha|\leq N}
      \frac{\partial_z^\alpha g(z)}{\alpha!}\xi^\alpha,
\]
\[
F_z^{\geq2}(\xi)
 =\sum_{2\leq|\alpha|\leq N+1}
      \frac{\partial_z^\alpha F(z)}{\alpha!}\xi^\alpha.
\]
Let
\[
H_F(z)=-F_{z_2}(z)\partial_{z_1}
                   +F_{z_1}(z)\partial_{z_2}.
\]
Taylor expansion and the product rule give the exact identity
\begin{equation}
\boxed{
j_N^+(H_Fg)
 =
H_F(z)\cdot\partial_z j_N^+g
 +\rho_N(F_z^{\geq2})j_N^+g.
}
\label{eq:hj-70}
\end{equation}
To verify it, write the full Taylor series \(Jg(z,\xi)\). Then
\[
J(H_Fg)=\{JF,Jg\}_\xi.
\]
Split \(JF\) into constant, linear, and degree-at-least-two terms. The constant contributes nothing. The linear term acts by
\(H_F(z)\cdot\partial_\xi Jg
 =H_F(z)\cdot\partial_zJg\).
Remove the constant output and truncate to degree \(N\). This is \eqref{eq:hj-70}.

Thus the geometric prolonged action is
\begin{equation}
\widetilde H_{F,N}
 =H_F(z)\cdot\partial_z+\rho_N(F_z^{\geq2}),
\label{eq:hj-71}
\end{equation}
not merely its matrix term. At a fixed point where \(dF=0\), the spatial vector field term vanishes on that fiber, recovering the origin-preserving jet module.
For example, take $F=z_1$, $g=z_2^2$, and $N=1$.
Then $F_z^{\geq2}=0$ and $j_1^+g=2z_2\xi_2$, whereas
\[
j_1^+(H_Fg)=2\xi_2
 =\partial_{z_2}(2z_2\xi_2).
\]
The matrix term is zero and the transport term supplies the entire action.

Arbitrary formal power series need not define holomorphic functions at moving spatial points. Equation \eqref{eq:hj-70} therefore assumes actual holomorphic functions, or a separately specified formal geometric setting.

A native geometric quantum-current comparison must construct the theory for the first-order spatial operators \eqref{eq:hj-71}, respect the relation defining holonomic jets, and compare its contractions with the matrix theory. The finite-dimensional internal current theorem does not supply those maps.

The next geometric construction must therefore retain both terms of
\eqref{eq:hj-71}. Its observable differential and contractions must
intertwine the holonomic jet map \eqref{eq:hj-70}; a matrix substitution
alone cannot provide that comparison.
